\documentclass[aspectratio=169]{beamer}
\usetheme{Boadilla}
\usecolortheme{dolphin}

\usepackage[utf8]{inputenc}
\usepackage[spanish]{babel}
\usepackage{booktabs}

\title{Guía Docente: Matriz de Diagnóstico}
\author{M.Sc. Ing. Bernardo Quiroga Turdera}
\institute{Circuitos Electrónicos III}

\begin{document}

\begin{frame}{Acciones Pedagógicas Recurrentes (Semanas 1-3)}
    \small
    \begin{table}[]
        \begin{tabular}{@{}p{2.5cm}p{3cm}p{7.5cm}@{}}
            \toprule
            \textbf{Deficiencia} & \textbf{Diagnóstico} & \textbf{Acción Pedagógica} \\ \midrule
            \textbf{Sec I: LVK/Thévenin} & Debilidad en redes & Reforzar Thévenin en divisores de tensión del BJT en Sem 2. \\ 
            \textbf{Sec II: Diodos} & Concepto idealizado & Usar modelo numérico no lineal de Shockley (Python) en Práctica 1. \\ 
            \textbf{Sec III: BJT} & Regiones de op. & Ejercicios en Sem 5 con rectas de carga interactivas en DC. \\ 
            \textbf{Sec IV: Frecuencia} & Dominio s vs t & Usar \texttt{scipy.signal} en Sem 4 para Bode dinámico, reduciendo abstracción. \\ 
            \textbf{Sec V: Python} & Poca experiencia & Asignar scripts base preconfigurados en Sem 1 para nivelar. \\ \bottomrule
        \end{tabular}
    \end{table}
\end{frame}

\end{document}